% !TEX program = xelatex
% !TEX root = main.tex
\documentclass[12pt,a4paper]{report}

\usepackage{ifxetex}
\ifxetex
  \usepackage{fontspec}
\else
  \usepackage[utf8]{inputenc}
  \usepackage[T1]{fontenc}
\fi
\usepackage[margin=1in]{geometry}
\usepackage{enumitem}
\usepackage{setspace}
\usepackage{titlesec}
\usepackage{fancyhdr}
\usepackage{graphicx}
\usepackage{float}
\usepackage{hyperref}
\usepackage{xurl}
\usepackage{longtable}
\usepackage{booktabs}
\usepackage{tabularx}
\usepackage{array}
\usepackage[section]{placeins}
\usepackage{xcolor}
\usepackage{ragged2e}
\usepackage{amssymb}
\ifxetex
\else
  \usepackage{microtype}
  \microtypesetup{expansion=false}
\fi

\setstretch{1.15}
\graphicspath{{figures/}}
\emergencystretch=2em
\setlength{\LTleft}{0pt}
\setlength{\LTright}{0pt}
\setlength{\tabcolsep}{4pt}
\setlength{\parindent}{0pt}
\setlength{\parskip}{0.45em}
\setlist[itemize]{leftmargin=1.4cm}
\setlist[enumerate]{leftmargin=1.4cm}
\sloppy
\titleformat{\chapter}[hang]
  {\normalfont\bfseries\Huge}
  {\thechapter}
  {0.8em}
  {}
\titleformat{\section}
  {\normalfont\bfseries\Large}
  {\thesection}
  {0.8em}
  {}
\titleformat{\subsection}
  {\normalfont\bfseries\large}
  {\thesubsection}
  {0.8em}
  {}
\titlespacing*{\chapter}{0pt}{-10pt}{22pt}
\titlespacing*{\section}{0pt}{18pt}{10pt}
\titlespacing*{\subsection}{0pt}{14pt}{8pt}
\let\cleardoublepage\clearpage
\renewcommand{\chaptername}{}
\newcolumntype{L}[1]{>{\RaggedRight\arraybackslash}p{#1}}
\newcolumntype{C}[1]{>{\centering\arraybackslash}p{#1}}
\newcommand{\codecell}[1]{{\ttfamily\small\path{#1}}}
\newcommand{\codeinline}[1]{{\ttfamily\path{#1}}}
\newcommand{\diagramimage}[1]{%
    \includegraphics[width=\textwidth,height=0.78\textheight,keepaspectratio]{#1}%
}

% --- TH Sarabun New Font Setup with Automatic Fallback for Overleaf Cloud ---
\ifxetex
  \IfFontExistsTF{TH Sarabun New}{
    \setmainfont{TH Sarabun New}[Scale=1.25, Script=Thai]
    \setsansfont{TH Sarabun New}[Scale=1.25, Script=Thai]
  }{
    \IfFontExistsTF{Laksaman}{
      \setmainfont{Laksaman}[Scale=1.2, Script=Thai]
      \setsansfont{Laksaman}[Scale=1.2, Script=Thai]
    }{
      \setmainfont{Garuda}[Scale=1.05, Script=Thai]
      \setsansfont{Garuda}[Scale=1.05, Script=Thai]
    }
  }
  \XeTeXlinebreaklocale "th"
  \XeTeXlinebreakskip = 0pt plus 1pt
\fi

\hypersetup{
    colorlinks=true,
    linkcolor=black,
    citecolor=black,
    urlcolor=black
}
\pagestyle{fancy}
\fancyhf{}
\renewcommand{\headrulewidth}{0pt}
\fancyfoot[R]{\thepage}

\begin{document}

\begin{titlepage}
    \centering
    \vspace*{2cm}
    {\LARGE \textbf{คู่มือการใช้งานระบบแพลตฟอร์ม "ก้าวใหม่"} \par}
    \vspace{0.5cm}
    {\Large (KaoMai Platform User Manual) \par}
    \vspace{0.3cm}
    {\large ฉบับผู้ใช้งานจริง | ฟอนต์ TH Sarabun \par}
    \vspace{2cm}
    {\large \textbf{เอกสารสำหรับ} \par}
    \vspace{0.5cm}
    เจ้าหน้าที่ศูนย์คนไร้บ้าน, ผู้จัดการรายกรณี (CM) \\
    และ ผู้จ้างงาน (Employer) \\
    \vspace{1.5cm}
    {\large \textbf{ขอบเขตการใช้งานในระบบ} \par}
    \vspace{0.5cm}
    การลงทะเบียนข้อมูลผู้เข้าร่วม, การจัดเก็บเอกสารประกอบ, \\
    การจับคู่งานอัจฉริยะ, แดชบอร์ด CM และการติดตามผลช่วงทำงาน \\
    \vfill
    {\large
    โครงการก้าวใหม่ (KaoMai Social Impact Initiative) \\
    ปีการศึกษา 2569
    }
\end{titlepage}

\pagenumbering{roman}

\chapter*{คำนำ}
\addcontentsline{toc}{chapter}{คำนำ}

เอกสารฉบับนี้คือ \textbf{"คู่มือการใช้งานระบบแพลตฟอร์มก้าวใหม่ (KaoMai User Manual)"} จัดทำขึ้นเพื่อให้ผู้ใช้งาน ได้แก่ เจ้าหน้าที่ศูนย์คนไร้บ้าน, ผู้จัดการรายกรณี (Case Manager - CM), และผู้จ้างงาน สามารถเรียนรู้และเข้าใจขั้นตอนการใช้งานส่วนต่าง ๆ ของระบบได้อย่างง่ายดาย

แพลตฟอร์มก้าวใหม่ถูกพัฒนาขึ้นเพื่อเป็นสื่อกลางในการเชื่อมโยงโอกาสงานระหว่างคนไร้บ้านที่พร้อมทำงานกับผู้จ้างงานที่ต้องการแรงงาน โดยมีฟังก์ชันอำนวยความสะดวกครบวงจร ตั้งแต่การบันทึกข้อมูลและเอกสารหลักฐาน การอนุมัติการจับคู่งาน ตลอดจนการติดตามผลการปรับตัวในสถานที่ทำงานอย่างเป็นระบบ

ผู้จัดทำหวังเป็นอย่างยิ่งว่าคู่มือฉบับนี้จะช่วยให้ผู้ใช้งานทุกท่านสามารถใช้งานระบบได้อย่างราบรื่นและเกิดประโยชน์สูงสุด

\clearpage
\tableofcontents
\clearpage
\pagenumbering{arabic}

\chapter{บทนำและภาพรวมการใช้งาน}

\section{จุดประสงค์ของแพลตฟอร์มก้าวใหม่}
แพลตฟอร์มก้าวใหม่ (KaoMai Web Application) ออกแบบมาเพื่อช่วยเหลือคนไร้บ้านให้เข้าถึงโอกาสงานที่สอดคล้องกับทักษะและความพร้อม โดยแบ่งบทบาทการใช้งานออกเป็น 2 กลุ่มหลัก:

\begin{itemize}
    \item \textbf{พื้นที่ศูนย์คนไร้บ้าน (Shelter Admin / Case Manager):} มีหน้าที่ดูแลข้อมูลผู้เข้าร่วม, มอบหมายผู้จัดการรายกรณี (CM), พิจารณาอนุมัติข้อเสนองานจากผู้จ้างงาน และติดตามผลการทำงานหลังได้รับการจ้างงาน
    \item \textbf{พื้นที่ผู้จ้างงาน (Employer):} มีหน้าที่สร้างประกาศรับสมัครงาน, คัดเลือกผู้สมัครจากคะแนนความเหมาะสม (Match Score \%) และติดต่อประสานงานเริ่มงาน
\end{itemize}

\section{เมนูหลักของระบบ}
\begin{enumerate}
    \item \textbf{สำหรับศูนย์คนไร้บ้าน:} เมนูผู้เข้าร่วม (\codeinline{/shelter/residents}), เมนูการดูแลรายกรณี (\codeinline{/shelter/case-management}), เมนูการจับคู่ (\codeinline{/shelter/matching}) และเมนูการติดตามงาน (\codeinline{/shelter/work-tracking})
    \item \textbf{สำหรับผู้จ้างงาน:} เมนูประกาศงานใหม่ (\codeinline{/employer/create-job}) และเมนูผู้สมัคร (\codeinline{/employer/matches})
\end{enumerate}

\chapter{คู่มือสำหรับศูนย์คนไร้บ้านและผู้จัดการรายกรณี}

\section{การลงทะเบียนและจัดการข้อมูลผู้เข้าร่วม}
\textbf{เมนู:} \texttt{ผู้เข้าร่วม} (\codeinline{/shelter/residents})

เจ้าหน้าที่ศูนย์พักพิงสามารถเพิ่มและดูแลข้อมูลคนไร้บ้านในความดูแลได้ตามขั้นตอนดังนี้:

\begin{enumerate}
    \item \textbf{ขั้นตอนการเพิ่มผู้เข้าร่วมใหม่:}
    \begin{itemize}
        \item กดปุ่ม \textbf{"+ เพิ่มผู้เข้าร่วม"} บริเวณมุมขวาบน
        \item \textbf{ข้อมูลพื้นฐาน:} อัปโหลดรูปถ่าย, กรอกชื่อที่ใช้แสดง, อายุ และเลือกเพศ
        \item \textbf{โรคประจำตัว:} กรอกข้อจำกัดสุขภาพหรือโรคประจำตัว (ข้อมูลนี้จะถูกเก็บเป็นความลับเฉพาะศูนย์พักพิง และไม่แสดงแก่นายจ้าง)
        \item \textbf{ผู้จัดการรายกรณี:} เลือกมอบหมาย Case Manager (CM) ผู้รับผิดชอบประจำเคส
        \item \textbf{สถานะบัตรประชาชน:} เลือกสถานะ (มีบัตรประชาชน / อยู่ระหว่างดำเนินการ / ต้องการความช่วยเหลือ)
        \item \textbf{รูปแบบงานที่ต้องการ:} เลือกระหว่าง \textbf{งานเต็มเวลา (Full-time)} หรือ \textbf{งานพาร์ทไทม์ (Part-time)}
        \item \textbf{วิธีรับค่าจ้างที่สะดวก:} เลือก \textbf{เงินสด} หรือ \textbf{โอนเข้าบัญชีธนาคาร}
        \item \textbf{ทักษะและเวลาที่พร้อม:} เลือกทักษะที่มี พร้อมกำหนดวันและช่วงเวลาที่สะดวกทำงาน
    \end{itemize}

    \item \textbf{ขั้นตอนการอัปโหลดเอกสารประกอบ (Supporting Documents):}
    \begin{itemize}
        \item ในส่วนเอกสารประกอบ กดปุ่ม \textbf{"อัปโหลดเอกสาร"}
        \item เลือกไฟล์จากคอมพิวเตอร์ (รองรับ PDF, JPG, PNG, WebP ขนาดไม่เกิน 5 MB สูงสุด 5 ไฟล์)
        \item เลือกประเภทเอกสาร เช่น \textit{วุฒิการศึกษา, ใบรับรองการฝึกอบรม, ประวัติการทำงาน หรือ เอกสารอื่น ๆ}
        \item สามารถกดปุ่ม \textbf{"เปิด"} เพื่อดาวน์โหลดดูไฟล์ หรือกด \textbf{"ลบ"} เพื่อยกเลิกเอกสารเดิมได้
    \end{itemize}
\end{enumerate}

\section{แดชบอร์ดการดูแลรายกรณี (Case Management Dashboard)}
\textbf{เมนู:} \texttt{การดูแลรายกรณี} (\codeinline{/shelter/case-management})

หน้าแดชบอร์ดนี้ช่วยให้ผู้จัดการรายกรณี (CM) และบริหารศูนย์ติดตามภาระงานได้อย่างเป็นระบบ:

\begin{itemize}
    \item \textbf{ตรวจสอบภาระงาน (Workload):} ดูสถิติจำนวนผู้เข้าร่วมที่ CM แต่ละคนกำลังดูแลอยู่
    \item \textbf{ตรวจสอบรายชื่อตกค้าง (Unassigned):} ระบบเตือนรายชื่อผู้เข้าร่วมที่ยังไม่มี CM ดูแล เพื่อให้กดมอบหมายผู้รับผิดชอบได้ทันที
    \item \textbf{ติดตามสถานะการประสานงาน:} ตรวจสอบผู้เข้าร่วมที่อยู่ระหว่างขั้นตอนรอคำตอบหรือรอเริ่มงาน
\end{itemize}

\section{การพิจารณาและอนุมัติการจับคู่งาน (Match Approval)}
\textbf{เมนู:} \texttt{การจับคู่} (\codeinline{/shelter/matching})

เมื่อผู้จ้างงานยื่นข้อเสนอมา รายการจะขึ้นในคิวการจับคู่:

\begin{enumerate}
    \item รายการข้อเสนอใหม่จะอยู่ในสถานะ \textbf{"ยื่นข้อเสนอแล้ว (Suggested)"}
    \item เจ้าหน้าที่ศูนย์หรือ CM เข้าไปสอบถามความสมัครใจและความพร้อมของผู้เข้าร่วม
    \item หากผู้เข้าร่วมพร้อมรับงาน ให้กดปุ่ม \textbf{"อนุมัติการจับคู่"} (ระบบจะเปิดเผยเบอร์โทรศัพท์ติดต่อศูนย์/CM แก่ผู้จ้างงานเพื่อเตรียมนัดหมาย)
    \item หากผู้เข้าร่วมไม่สะดวกรับงาน ให้กดปุ่ม \textbf{"ไม่อนุมัติ / ผู้สมัครไม่สะดวก"}
\end{enumerate}

\section{ระบบติดตามการทำงานและการยุติการทำงาน (Work Tracking)}
\textbf{เมนู:} \texttt{การติดตามงาน} (\codeinline{/shelter/work-tracking})

หลังจากคนไร้บ้านเริ่มทำงานแล้ว ให้กดปุ่ม \textbf{"+ เริ่มติดตามงาน"} เพื่อเข้าสู่ระบบดูแลต่อเนื่อง:

\begin{enumerate}
    \item \textbf{การตั้งค่ารอบติดตามงาน:}
    \begin{itemize}
        \item \textbf{สัปดาห์แรก (1st Week):} ติดตามและบันทึกผลการทำงานรายวัน
        \item \textbf{รอบติดตามถัดไป:} เลือกระยะเวลาเป็น \textbf{ทุก 2 สัปดาห์ (Fortnightly)} หรือ \textbf{ทุกเดือน (Monthly)} โดยระบบจะคำนวณวันนัดติดตามครั้งถัดไปให้อัตโนมัติ
    \end{itemize}

    \item \textbf{การลงบันทึกการติดตาม (Check-In):}
    \begin{itemize}
        \item กดปุ่ม \textbf{"บันทึกติดตาม"}
        \item ระบุการมาทำงาน: \textit{มาทำงานปกติ / มาสาย / ขาดงาน}
        \item ระบุระดับการปรับตัว: \textit{ปรับตัวได้ดี / ต้องการการสนับสนุน / เร่งด่วน}
        \item กรอกความคิดเห็นจากผู้เข้าร่วม ความคิดเห็นจากนายจ้าง และบันทึกภายในสำหรับทีมศูนย์
    \end{itemize}

    \item \textbf{การยุติการทำงาน (End Employment):}
    \begin{itemize}
        \item เมื่อสิ้นสุดสัญญาหรือหยุดทำงาน ให้กดปุ่ม \textbf{"ยุติการทำงาน"}
        \item เลือกเหตุผลการยุติ: \textit{สิ้นสุดสัญญาตามกำหนด / ผู้เข้าร่วมลาออก / นายจ้างยุติการจ้าง / เหตุผลสุขภาพส่วนตัว / ขาดการติดต่อ / อื่น ๆ}
        \item กรอกสรุปโน้ตปิดเคส
        \item หากต้องการส่งผู้เข้าร่วมกลับเข้าสู่ระบบหางานใหม่ ให้เลือกติ๊กถูกที่ \textbf{"กลับเข้าสู่กระบวนการหางาน"} เพื่อส่งรายชื่อกลับมาพร้อมรับงานใหม่อีกครั้ง
    \end{itemize}
\end{enumerate}

\chapter{คู่มือสำหรับผู้จ้างงาน}

\section{การสร้างและลงประกาศงานใหม่ (Create Job Posting)}
\textbf{เมนู:} \texttt{ประกาศงานใหม่} (\codeinline{/employer/create-job})

ผู้จ้างงานสามารถลงประกาศรับสมัครงานได้ตามขั้นตอนดังนี้:

\begin{enumerate}
    \item กรอกชื่อตำแหน่งงาน (เช่น ช่างไม้, พนักงานทำความสะอาด, จัดเรียงสินค้า)
    \item รายละเอียดงาน สถานที่ปฏิบัติงาน และกำหนดค่าจ้างรายวัน (บาท/วัน)
    \item เลือกทักษะที่จำเป็นสำหรับตำแหน่งงาน แล้วกด \textbf{"ลงประกาศงาน"}
\end{enumerate}

\section{การคัดเลือกผู้สมัครและการยื่นข้อเสนองาน (Smart Matching)}
\textbf{เมนู:} \texttt{ผู้สมัคร} (\codeinline{/employer/matches})

หน้าคัดเลือกผู้สมัครถูกแบ่งออกเป็นแท็บเพื่อความสะดวกในการติดตาม:

\begin{enumerate}
    \item \textbf{แท็บ "ผู้สมัครแนะนำ":}
    \begin{itemize}
        \item ระบบจะแสดงรายชื่อคนไร้บ้านที่ผ่านการคัดกรอง เรียงลำดับตามคะแนนความเหมาะสม (Match Score \%) และระยะทางโดยประมาณจากศูนย์พักพิง
        \item แสดงแท็กทักษะที่ตรงกับงาน ($\checkmark$) สถานะบัตรประชาชน และรูปโปรไฟล์
        \item กดปุ่ม \textbf{"ดูรายละเอียดโปรไฟล์"} เพื่อดูประวัติอย่างละเอียด
        \item หากสนใจผู้สมัครคนใด ให้กดปุ่ม \textbf{"ยื่นข้อเสนองาน"} เพื่อส่งข้อเสนอไปยังศูนย์คนไร้บ้าน
    \end{itemize}

    \item \textbf{แท็บ "ยื่นเสนอแล้ว/รอตอบ":} สำหรับติดตามรายการข้อเสนองานที่อยู่ระหว่างรอศูนย์คนไร้บ้านสอบถามผู้สมัคร
    \item \textbf{แท็บ "อนุมัติแล้ว/พร้อมเริ่มงาน":} เมื่อศูนย์คนไร้บ้านกดอนุมัติแล้ว นายจ้างจะสามารถกดดูเบอร์โทรศัพท์ติดต่อศูนย์พักพิงและ Case Manager เพื่อเตรียมนัดหมายเริ่มงานได้ทันที
\end{enumerate}

\end{document}
